%Daten zur Institution

%\input{Dozentdaten}

%\renewcommand{\fachbereich}{Fachbereich}

%\renewcommand{\dozent}{Prof. Dr. . }

%Klausurdaten

\renewcommand{\klausurgebiet}{ }

\renewcommand{\klausurtyp}{ }

\renewcommand{\klausurdatum}{. 20}

\klausurvorspann {\fachbereich} {\klausurdatum} {\dozent} {\klausurgebiet} {\klausurtyp}

%Daten für folgende Punktetabelle

\renewcommand{\aeins}{  3  }

\renewcommand{\azwei}{  3   }

\renewcommand{\adrei}{  6  }

\renewcommand{\avier}{  5  }

\renewcommand{\afuenf}{  5  }

\renewcommand{\asechs}{  2  }

\renewcommand{\asieben}{  5  }

\renewcommand{\aacht}{  5  }

\renewcommand{\aneun}{  7  }

\renewcommand{\azehn}{  3  }

\renewcommand{\aelf}{  2  }

\renewcommand{\azwoelf}{  5  }

\renewcommand{\adreizehn}{  9   }

\renewcommand{\avierzehn}{  5  }

\renewcommand{\afuenfzehn}{  65  }

\renewcommand{\asechzehn}{    }

\renewcommand{\asiebzehn}{    }

\renewcommand{\aachtzehn}{    }

\renewcommand{\aneunzehn}{    }

\renewcommand{\azwanzig}{    }

\renewcommand{\aeinundzwanzig}{    }

\renewcommand{\azweiundzwanzig}{    }

\renewcommand{\adreiundzwanzig}{    }

\renewcommand{\avierundzwanzig}{    }

\renewcommand{\afuenfundzwanzig}{    }

\renewcommand{\asechsundzwanzig}{    }

\punktetabellevierzehn

\klausurnote

\newpage

\setcounter{section}{0}

\inputaufgabegibtloesung
{3}
{

Definiere die folgenden
\zusatzklammer {kursiv gedruckten} {} {} Begriffe.
\aufzaehlungsechs{Der
\stichwort {Krümmungskreis} {}
zu einer zweifach
\definitionsverweis {stetig differenzierbare}{}{}
\definitionsverweis {bogenparametrisierte}{}{}
\definitionsverweis {Kurve}{}{}
Es sei
\maabbdisp {\gamma} {I} {\R^2
} {}
in einem Punkt

\mavergleichskette
{\vergleichskette
{ t_0
}
{ \in }{ I
}
{  }{
}
{    }{
}
{   }{
}
}
{}{}{.}

}{Eine \stichwort {abgeschlossene Untermannigfaltigkeit} {}
\mathl{M \subseteq N}{} einer differenzierbaren Mannigfaltigkeit $N$.

}{Der \stichwort {Tangentialraum} {} in einem Punkt
\mathl{P\in M}{} einer differenzierbaren Mannigfaltigkeit $M$.

}{Eine
\stichwort {orientierte} {}
\definitionsverweis {differenzierbare Mannigfaltigkeit}{}{}
$M$.

}{Der
\stichwort {Rand} {}
einer
\definitionsverweis {Mannigfaltigkeit mit Rand}{}{.}

}{Die
\stichwort {Christoffelsymbole} {}

\mathl{\Gamma_{ij}^k}{} für den Levi-Civita-Zusammenhang zu einer riemannschen Struktur  $\left(  g_{ij}  \right)$ auf einer offenen Menge

\mavergleichskette
{\vergleichskette
{ U
}
{ \subseteq }{ \R^n
}
{  }{
}
{    }{
}
{   }{
}
}
{}{}{.}
}

}
{} {}

\inputaufgabegibtloesung
{3}
{

Formuliere die folgenden Sätze.
\aufzaehlungdrei{Der Satz über das Bild der
\definitionsverweis {Gauß-Abbildung}{}{}
eine Hyperfläche.}{Die Formel für den Flächeninhalt einer Rotationsfläche zu einer
\definitionsverweis {differenzierbaren Kurve}{}{}
\maabbeledisp {\gamma} {]a,b[} {\R^2
} { t } {(x(t),y(t))
} {,}
mit

\mavergleichskette
{\vergleichskette
{  y(t)
}
{ > }{  0
}
{  }{
}
{    }{
}
{   }{
}
}
{}{}{.}}{Die
\stichwort {Produktregel} {}
für die äußere Ableitung.}

}
{} {}

\inputaufgabegibtloesung
{6 (1+2+3)}
{

Es sei $Y$ die Hyperfläche, die durch die Bedingung

\mavergleichskettedisp
{\vergleichskette
{ 2x^2 +3y^2+5z^2
}
{ =} { 10
}
{ } {
}
{ } {
}
{ } {
}
}
{}{}{}
im $\R^3$ gegeben ist. Es sei

\mavergleichskette
{\vergleichskette
{ P
}
{ = }{ \begin{pmatrix}  -1 \\ -1\\  1   \end{pmatrix}
}
{ \in }{ Y
}
{    }{
}
{   }{
}
}
{}{}{}
und

\mavergleichskette
{\vergleichskette
{ w
}
{ = }{ \begin{pmatrix}  -3 \\ 2 \\  4   \end{pmatrix}
}
{ \in }{ \R^3
}
{    }{
}
{   }{
}
}
{}{}{.}
\aufzaehlungdrei{Bestimme den
\definitionsverweis {Tangentialraum}{}{}

\mathl{T_PY}{.}
}{Bestimme die
\definitionsverweis {orthogonale Zerlegung}{}{}
$w$ in die tangentiale und in die orthogonale Komponente zum Punkt $P$.
}{Realisiere die tangentiale Komponente von $w$ in $T_PY$ durch eine differenzierbare Kurve, die ganz auf $Y$  verläuft.
}

}
{} {}

\inputaufgabegibtloesung
{5}
{

Es sei
\maabbeledisp {\gamma} {\R} { \R^2
} { t } { \begin{pmatrix}  \cos \left(  \omega t \right)  \\ \sin  \left(  \omega t \right)    \end{pmatrix}
} {,}
mit einem

\mavergleichskette
{\vergleichskette
{ \omega
}
{ > }{ 0
}
{  }{
}
{    }{
}
{   }{
}
}
{}{}{.}
Man zeige, dass es keine andere Kreisbewegung mit konstanter Geschwindigkeitsnorm gibt, die mit $\gamma$ für

\mavergleichskette
{\vergleichskette
{ t
}
{ = }{ 0
}
{  }{
}
{    }{
}
{   }{
}
}
{}{}{}
bis zur zweiten Ableitung übereinstimmt.

}
{} {}

\inputaufgabegibtloesung
{5}
{

Beweise den Satz über die Krümmung auf einer implizit gegebenen ebenen Kurve.

}
{} {}

\inputaufgabegibtloesung
{2}
{

Man gebe ein Beispiel für differenzierbare Mannigfaltigkeiten
\mathkor {} {M} {und} {N} {}
mit

\mavergleichskette
{\vergleichskette
{ \operatorname{ dim}_{  } ^{  }  \left(   M   \right), \, \operatorname{ dim}_{  } ^{  }  \left(   N   \right)
}
{ \geq }{ 1
}
{  }{
}
{    }{
}
{   }{
}
}
{}{}{}
und differenzierbare Abbildungen
\maabb {\varphi} { M } { N
} {}
und
\maabb {\psi} { N } { M
} {}
derart, dass

\mavergleichskette
{\vergleichskette
{ \psi \circ \varphi
}
{ = }{
\operatorname{Id}_{ M }
}
{  }{
}
{    }{
}
{   }{
}
}
{}{}{}
und

\mavergleichskette
{\vergleichskette
{\varphi \circ  \psi
}
{ \neq }{
\operatorname{Id}_{ N }
}
{  }{
}
{    }{
}
{   }{
}
}
{}{}{}
gilt.

}
{} {}

\inputaufgabegibtloesung
{5}
{

Beweise die Aussage, dass die tangentiale Äquivalenz von Wegen auf einer Mannigfaltigkeit in einem Punkt $P$ mit einer beliebigen Karte überprüft werden kann.

}
{} {}

\inputaufgabegibtloesung
{5}
{

Zeige, dass die Tangentialabbildung $T(\varphi)$ zu
\maabbeledisp {\varphi} {\R^1} {S^1
} { t } {( \cos  t, \sin  t  )
} {,}
surjektiv ist.

}
{} {}

\inputaufgabegibtloesung
{7 (3+4)}
{

Wir betrachten die differenzierbaren Abbildungen
\maabbeledisp {\gamma} { [1,c] } {\R^2
} { t } {(t,t^{3})
} {,}
\zusatzklammer {mit \mathlk{c \geq 1}{}} {} {}
und
\maabbeledisp {\varphi} {\R_+ \times \R_+ } {\R^3
} {(u,v)} {(u^3,u^2+v^2,u^{-1}v^{-1} )
} {,}
und die Differentialform

\mavergleichskettedisp
{\vergleichskette
{ \omega
}
{ =} { (x-y)dx-z^2dy+dz
}
{ } {
}
{ } {
}
{ } {
}
}
{}{}{}
auf dem $\R^3$.

a) Berechne die zurückgezogene Differentialform
\mathl{\varphi^* \omega}{} auf dem $\R_+ \times \R_+$.

b) Berechne das Wegintegral zur Differentialform
\mathl{\varphi^* \omega}{} zum Weg $\gamma$ in Abhängigkeit von $c$.

}
{} {}

\inputaufgabegibtloesung
{3}
{

Es sei
\maabbdisp {\psi} { L } { M
} {}
eine
\definitionsverweis {differenzierbare Abbildung}{}{}
zwischen
\definitionsverweis {differenzierbaren Mannigfaltigkeiten}{}{,}
es sei
\maabbdisp {f} { M } {\R
} {}
eine Funktion auf $M$ und $\omega$ eine
$k$-\definitionsverweis {Form}{}{}
auf $M$. Zeige

\mavergleichskettedisp
{\vergleichskette
{ \psi^*(f \omega)
}
{ =} { \psi^*(f)  \psi^* \omega
}
{ } {
}
{ } {
}
{ } {
}
}
{}{}{.}

}
{} {}

\inputaufgabe
{2}
{

Beschreibe möglichst viele
\definitionsverweis {zusammenhängende}{}{}
eindimensionale
\definitionsverweis {Mannigfaltigkeiten mit Rand}{}{.}

}
{} {}

\inputaufgabegibtloesung
{5}
{

Es sei
\maabbdisp {f} { [a,b] } { \R_{+}
} {}
eine
\definitionsverweis {stetig differenzierbare Funktion}{}{.}
Wir fassen den
\definitionsverweis {Subgraphen}{}{}
als eine
\definitionsverweis {Mannigfaltigkeit mit Rand}{}{}
auf, wobei der Rand aus dem Graphen, dem Grundintervall und den beiden Seitenkanten, aber ohne die vier Eckpunkte besteht. Bestätige den Satz von Stokes direkt für die
\definitionsverweis {Differentialform}{}{}

\mavergleichskette
{\vergleichskette
{ \omega
}
{ = }{ xdy
}
{  }{
}
{    }{
}
{   }{
}
}
{}{}{.}

}
{} {}

\inputaufgabegibtloesung
{9}
{

Beweise den Brouwerschen Fixpunktsatz.

}
{} {}

\inputaufgabegibtloesung
{5 (3+1+1)}
{

Wir betrachten den Zylinder

\mavergleichskette
{\vergleichskette
{ Z
}
{ = }{ S^1 \times \R
}
{ \subseteq }{ \R^2 \times \R
}
{    }{
}
{   }{
}
}
{}{}{}
mit dem
\definitionsverweis {trivialen Zusammenhang}{}{}
auf seinem
\definitionsverweis {Tangentialbündel}{}{.}
\aufzaehlungdrei{Bestimme in einer geeigneten isometrischen Karte die
\definitionsverweis {geodätische Differentialgleichung}{}{}
für diejenige geodätische Kurve $\psi$  auf $M$,  die

\mavergleichskettedisp
{\vergleichskette
{ \psi(0)
}
{ =} { (1,0,3)
}
{ } {
}
{ } {
}
{ } {
}
}
{}{}{}
und

\mavergleichskettedisp
{\vergleichskette
{ \psi'(0)
}
{ =} { (0,1,-4)
}
{ } {
}
{ } {
}
{ } {
}
}
{}{}{}
erfüllt.
}{Löse die Differentialgleichung aus (1) in der Karte.
}{Finde eine geodätische Kurve in $Z$, die die Bedingungen aus (1) erfüllt.
}

}
{} {}

 [[Kategorie:Latexseite]]
